\section{Conclusion}
\label{sec:conclusion}

Honeypots have spent two decades recording attacker behavior without
reasoning about it; generative models opened a worse one: unverified
improvisation as the product. This paper argued that the right unit of
design is the \emph{loop}, and the evaluation supports that division of
labor: an echo guard yields echo-corrected lower bounds of 0.479/0.243;
retrieval grounding, once removed, costs $+0.154$ and every TTP; and
under pressure the fiction graph holds---zero
contradictions in 45 re-queries, zero unverified serves in 218 live
queries, every ablation costing its property---and deception quality
holds $3.86/5$; the verifier's misses are published (540 of 572 caught,
17 of 206 benign refused, both injection cases missed). At the top of
the ladder the result is a subtraction: kind-awareness fixed 12/13
real-log assessments and lifted bait exposure 125\% in five replays but
added no dwell and left the path; the Evidence Judge stayed disabled
after a prospective null.

Ambient traffic is scanner-grade only, so engagement evidence rests on
self-play and two past attackers; every labeling lane is
machine-generated---the panel reports $\kappa=0.81$ within one family,
its oracle agreeing with the majority on only $50.8\%$ of cases---no
label carries human authority. The gate's counterweights are measured:
43 fail-closed refusals on natural traffic, zero leaked credentials,
one-in-twelve out-of-sample recall, one confirmed deception-plane yield.

Adoption follows the ladder's rule: canary registry first, floor
authoritative, further rungs only through pre-registered gates. The
contribution is a budget: three rungs kept, three removed, and the
discipline telling an operator which is which.
